\section{Conclusion}\label{sec:conclusion}

We studied how reducing the search space helps solve the single-source capacitated facility
location problem and how much of that help comes from learning. On our data, safe pruning
required keeping about 80\% of the candidate facilities, and decomposition without coordination
produced infeasible unions. Two mechanisms worked: an adaptive expansion over ranked facilities
with a reduced-cost optimality certificate, and a cluster-oriented LNS that re-optimizes one
subproblem at a time over residual capacities. In validation pilots both were clearly better than
the solver on the full model. For the learned components the pilots found no evidence of a benefit: the point estimates
favored the linear-relaxation ranking over the GNN, and learned subproblem selection did not
separate from rotation.
More than half of the subproblems selected by the neighborhood search had already been tried
without success from the same local state, which suggests, without proving it, that the search
is limited by the reach of its neighborhoods more than by their order.

The pre-registered closed test, on 139 instances that no pilot had used, sharpened these
conclusions. Restricting the facilities by the LP ranking more than halved the primal integral
of the full solver. Replacing the LP ranking by the GNN made the expansion significantly worse,
increasingly so away from the training distribution, to the point of matching the full solver on
larger instances. Kernel search, a learning-free method from 2014, was at the level of the nested
expansion. The neighborhood-search phase did not change the primal integral, lowered the final
deviation by about half a percentage point at the threshold of significance, and was the only
component that came close to the best known solution on real instances with 600 and 850
customers (four instances, reported descriptively).

The practical recommendation is to report, for any learned search-space reduction, the
learning-free method with the same function and the same budget. Future work includes
neighborhoods that open and close several facilities at once, an online estimate of the opening
distribution from the solutions visited during the search in place of an offline-trained ranker,
larger real instances where the full model no longer fits and decomposition becomes mandatory,
and extending the coordinated decomposition to location-routing, where our reconstruction
indicates that the classical continuous approximation is the level to beat.
